\begin{lemma}[A common joint-flux measurement on double-exchange outputs]%
    \label{thm:peps_torus_three_plaquette_output_measurement}
    \lean{TNLean.PEPS.IsGIsometric.exists_torusThreePlaquetteOutput_physicalMeasurement}
    \leanok
    \uses{thm:peps_three_plaquette_geometry,
        thm:peps_torus_three_plaquette_flux_holonomy,
        thm:peps_torus_three_plaquette_double_exchange_holonomy,
        thm:peps_closed_walk_class_physical_measurement,
        thm:peps_physical_cut_column_action}
    Let $G$ be finite and let the native four-leg tensor be $G$-isometric
    for the regular action. On a torus of horizontal period at least five
    and vertical period at least three, fix any translated eight-site strip
    $R_v$. For directed bond values $t$, write
    \begin{align}
        (a,b,c) & =(t_0^{-1}t_1,t_1^{-1}t_2,t_2^{-1}), \\
        k(t)    & =aba^{-1}c.
    \end{align}
    There is one complete projective measurement
    $(Q_C)_{C\in\operatorname{Conj}(G)\cup\{\bot\}}$ on the original spins
    of $R_v$, chosen before every $t$, common exterior and crossing
    assignment $u$, and boundary configuration $\theta$.
    Each $Q_C$ is positive and self-adjoint, with
    \begin{align}
        Q_CQ_D & =\delta_{C,D}Q_C, & \sum_C Q_C & =I.
    \end{align}
    On the output bond values
    $t'=(t_0,t_2t_0^{-1}t_1t_2^{-1}t_0,t_2)$, this family satisfies
    \begin{align}
        Q_C\Psi_{R_v}(t',u,\theta)
                   & =\mathbf1_{C=[k(t)]}\Psi_{R_v}(t',u,\theta), \\
        Q_CM(t',u) & =\mathbf1_{C=[k(t)]}M(t',u),
    \end{align}
    where $M$ is the global physical cut matrix.
    The unused outcome $\bot$ annihilates every such column and cut.
    All positions and both seams are included, with no auxiliary tree,
    holonomy, Gram or state hypothesis.
    This is the joint-measurement step for the constructed auxiliary
    double exchange in \cite[braiding passage]{Schuch2010PEPS}, using
    \cite[Theorem 6.15]{Schuch2010PEPS}. Identification with a prescribed
    physical string braid, energies and parent-Hamiltonian membership
    remain separate.
\end{lemma}

\begin{proof}\leanok
    The outer walk surrounding the middle and right plaquettes lies
    in $R_v$, whose induced graph is connected. Inducing this closed walk
    in the region and applying the common closed-walk measurement theorem,
    internal restriction removes every common exterior and crossing
    operator from its holonomy, while the output formula gives $k(t)$.
    The uniform boundary-column action then passes to the global cut by its
    open-region contraction identity.
\end{proof}

\begin{lemma}[The same physical double exchange followed by joint measurement]%
    \label{thm:peps_torus_three_plaquette_measured_double_exchange}
    \lean{TNLean.PEPS.IsGIsometric.exists_torusThreePlaquetteMeasuredDoubleExchange}
    \leanok
    \uses{thm:peps_torus_three_plaquette_output_measurement,
        thm:peps_torus_three_plaquette_double_exchange_physical}
    Under the preceding hypotheses, the measurement $Q$ can be chosen
    together with one unitary $W$ implementing the double exchange.
    Both operators are chosen before every $t,u,\theta$. The same $W$ satisfies
    \begin{align}
        W\Psi_{R_v}(t,\theta)         & =\Psi_{R_v}(t',\theta),                    \\
        (W\otimes I_{R_v^c})\Psi(t,u) & =\Psi(t',u),                               \\
        Q_CW\Psi_{R_v}(t,\theta)
                                      & =\mathbf1_{C=[k(t)]}W\Psi_{R_v}(t,\theta), \\
        Q_CWM(t,u)                    & =\mathbf1_{C=[k(t)]}WM(t,u).
    \end{align}
    Its identity extension is unitary. In particular, input plaquette
    holonomies $(g,h,h^{-1})$ give the joint eigenspace
    $[ghg^{-1}h^{-1}]$ after this constructed operation.
    The scope is the same auxiliary double exchange, without a
    prescribed string-crossing or parent-Hamiltonian assertion.
\end{lemma}

\begin{proof}\leanok
    Choosing $W$ from the double-exchange theorem and $Q$ from the
    output measurement theorem, their uniform quantifiers allow both
    choices before the bond values and boundary data. Substituting the local
    column identity into the measurement action gives the column equation, and
    evaluating the global identity on every combined cut configuration yields
    $WM(t,u)=M(t',u)$, which substitutes into the output cut measurement.
\end{proof}
